\documentclass[11pt,a4paper]{article}

% Packages
\usepackage[utf8]{inputenc}
\usepackage[T1]{fontenc}
\usepackage{geometry}
\usepackage{titlesec}
\usepackage{enumitem}
\usepackage{hyperref}
\usepackage{xcolor}
\usepackage{tabularx}
\usepackage{booktabs}
\usepackage{fancyhdr}
\usepackage{lastpage}
\usepackage{listings}
\usepackage{multicol}

% Page Geometry
\geometry{top=1in, bottom=1in, left=1in, right=1in}

% Colors
\definecolor{primary}{RGB}{23, 75, 75}
\definecolor{secondary}{RGB}{34, 91, 88}
\definecolor{lightbg}{RGB}{234, 242, 238}
\definecolor{codegray}{rgb}{0.5,0.5,0.5}

% Code Block Styling
\lstset{
    backgroundcolor=\color{lightbg},
    basicstyle=\ttfamily\small,
    breaklines=true,
    captionpos=b,
    commentstyle=\color{codegray},
    keywordstyle=\color{primary},
    stringstyle=\color{secondary},
    frame=single,
    rulecolor=\color{primary}
}

% Hyperlink Setup
\hypersetup{colorlinks=true, linkcolor=primary, filecolor=primary, urlcolor=primary}

% Header and Footer
\pagestyle{fancy}
\fancyhf{}
\fancyhead[L]{\textcolor{secondary}{\textbf{LossLens}}}
\fancyhead[R]{\textcolor{secondary}{Detailed Technical \& Feature Documentation}}
\fancyfoot[C]{\textcolor{secondary}{Page \thepage\ of \pageref{LastPage}}}
\renewcommand{\headrulewidth}{0.4pt}

% Section Formatting
\titleformat{\section}{\Large\bfseries\color{primary}}{\thesection}{1em}{}[\titlerule]
\titleformat{\subsection}{\large\bfseries\color{secondary}}{\thesubsection}{1em}{}
\titleformat{\subsubsection}{\normalsize\bfseries\color{secondary}}{\thesubsubsection}{1em}{}

\begin{document}

\begin{titlepage}
    \centering
    \vspace*{2cm}
    {\Huge \textbf{\textcolor{primary}{LossLens}} \par}
    \vspace{0.5cm}
    {\Large \textbf{Detailed Technical \& Feature Documentation} \par}
    \vspace{1.5cm}
    {\large \textbf{Version:} 2.0 (Pre-Pilot Technical Blueprint) \par}
    {\large \textbf{Date:} November 2026 \par}
    {\large \textbf{Architecture:} React (Frontend), Supabase (Backend), Capacitor (Mobile wrapper) \par}
    \vfill
\end{titlepage}

\newpage
\tableofcontents
\newpage

\section{Executive Summary}
This document serves as the exhaustive technical and operational blueprint for the LossLens ecosystem. It is designed for software developers, clinical directors, insurance underwriters, and competition judges. It details the precise state machines, data flows, API integrations, and security protocols required to launch a fraud-proof, privacy-centric gambling intervention platform.
% INSERT THIS SECTION AFTER THE EXECUTIVE SUMMARY IN PART 1

\section{Core Platform Features (Existing)}
LossLens is already a functioning digital recovery platform. The following features are live and form the foundation of the ecosystem.

\subsection{The AI Companion (Oloo)}
\begin{itemize}
    \item \textbf{24/7 Conversational Support:} Oloo is a generative AI companion designed to provide non-judgmental, immediate support to users experiencing gambling urges or emotional distress.
    \item \textbf{CBT Integration:} Oloo is instructed via system prompts to use Cognitive Behavioral Therapy (CBT) techniques, such as reframing negative thoughts and guiding the user through urge-surfing exercises.
    \item \textbf{Safety Guardrails:} Oloo is programmed to recognize crisis language (e.g., self-harm) and immediately halt the generative chat to present emergency hotlines (Befrienders Kenya) and grounding exercises.
\end{itemize}

\subsection{CBT-Inspired Learning Modules}
\begin{itemize}
    \item \textbf{Structured Curriculum:} A library of self-guided learning modules that teach users the mechanics of gambling addiction (e.g., the sunk cost fallacy, variable ratio reinforcement).
    \item \textbf{Interactive Exercises:} Users complete worksheets and reflections within the app to identify their personal triggers and build coping strategies.
    \item \textbf{Progress Tracking:} The app tracks module completion and feeds this data into the overall user profile to measure engagement.
\end{itemize}

\subsection{Personal Reflection \& Diary Tools}
\begin{itemize}
    \item \textbf{Trigger Journaling:} A private diary space for users to log their feelings, identify what triggered an urge, and reflect on their recovery journey.
    \item \textbf{Daily Mood Tracking:} A quick-check UI where users log their mood and urge intensity. (Note: This data feeds the actuarial Markov model for risk prediction).
\end{itemize}

\subsection{Peer \& Community Support}
\begin{itemize}
    \item \textbf{Community Rooms:} Secure, moderated discussion forums where users can share their struggles and successes anonymously.
    \item \textbf{Moderation Protocol:} Community guidelines strictly prohibit predatory behavior, gambling promotion, or sharing of betting tips. 
\end{itemize}

\subsection{Recovery Resources \& Professional Pathways}
\begin{itemize}
    \item \textbf{Resource Library:} A curated list of articles, videos, and tools for financial planning and mental wellness.
    \item \textbf{Professional Help Directory:} A pathway that directs users to GAMHELP Kenya for long-term, in-person rehabilitation and support groups.
\end{itemize}

\subsection{System Utility \& PWA Features}
\begin{itemize}
    \item \textbf{Progressive Web App (PWA):} LossLens is installable on the home screen of a user's phone without needing an app store.
    \item \textbf{Service Worker Architecture:} The app caches static assets (like CBT modules) so that users can access recovery tools even when they are offline or have zero data.
    \item \textbf{Push Notifications:} Automated check-in reminders and motivational prompts sent directly to the user's device.
\end{itemize}
\section{The User Application (Mobile/Web App)}
The User App is a Progressive Web App (PWA) built in React, wrapped via Capacitor for native Android/iOS deployment. It functions primarily as a behavioral data collection and intervention tool.

\subsection{State Machine and User Flow}
The app operates on a strict state machine driven by the user's PGSI score and daily behavioral data.
\begin{enumerate}
    \item \textbf{State 0: Unregistered} - User can access public resources, read CBT articles, and dial crisis lines.
    \item \textbf{State 1: Onboarded (Free Tier)} - User completed PGSI. App tracks data. Markov model calculates risk in background.
    \item \textbf{State 2: Elevated Risk} - Markov model detects a transition toward the "Lapsed" state. App suggests upgrading to Premium.
    \item \textbf{State 3: Premium (Active)} - User pays KES 250/month. 30-day waiting period for clinical coverage begins.
    \item \textbf{State 4: Clinical Crisis} - Risk score crosses critical threshold OR user triggers Panic Button. Digital Voucher generated.
\end{enumerate}

\subsection{Onboarding \& PGSI Stratification}
\begin{itemize}
    \item \textbf{PGSI Form:} 9 questions. Scores range from 0 to 27.
    \item \textbf{Risk Bands:}
    \begin{itemize}
        \item 0-2: Low Risk. UI shows general wellness tools.
        \item 3-7: Moderate Risk. UI highlights CBT modules and daily check-ins.
        \item 8+: High Risk. UI unlocks "Red Zone" protocol instantly. Premium upsell is aggressive.
    \end{itemize}
    \item \textbf{Consent Flow:} Two separate checkboxes. (1) Terms of Service. (2) Privacy \& Data Processing (aligned with Kenya Data Protection Act 2019). User must explicitly consent to AI chat processing via third-party LLMs (Groq).
\end{itemize}

\subsection{Daily Check-in \& AI Companion (Oloo)}
\begin{itemize}
    \item \textbf{Daily Check-in Data Points:}
    \begin{itemize}
        \item Mood score (1-5)
        \item Urge intensity (1-10)
        \item Gambled today? (Y/N)
        \item Optional trigger tags (e.g., "Payday", "Stress", "Boredom")
    \end{itemize}
    \item \textbf{AI Chat (Oloo):}
    \begin{itemize}
        \item Built on Groq API for ultra-low latency responses.
        \item System Prompt: Guardrails prevent Oloo from diagnosing. Oloo redirects crisis language to hotlines.
        \item Context Window: Oloo is fed the user's last 3 daily check-ins to provide contextual support (e.g., "I see your urges were high yesterday.").
    \end{itemize}
    \item \textbf{Data Pipeline:} Every check-in and chat summary is written to the Supabase \texttt{behavioral\_logs} table. This data is fed into the Markov model nightly via a cron job.
\end{itemize}

\subsection{The LossLens Shield (Native Integration)}
This is the app's core prevention mechanism. It relies on native OS permissions (Android Accessibility Service / iOS Screen Time API).
\begin{itemize}
    \item \textbf{Activation:} User toggles "Enable Shield" in settings and grants OS permissions.
    \item \textbf{Trigger Logic:} The Shield activates dynamically if:
    \begin{enumerate}
        \item The user logs an urge intensity of 8 or higher.
        \item The user opens a known betting app (Betika, SportPesa, Odibets) between 6 PM and 4 AM (peak gambling hours).
        \item The Markov model flags a 30\% probability of relapse.
    \end{enumerate}
    \item \textbf{The Overlay:} When a betting app is opened, a full-screen LossLens overlay interrupts the launch.
    \item \textbf{The Friction Unlock:} The user must complete a 2-minute guided breathing exercise and type a 1-sentence reflection on why they shouldn't bet. Only then does the overlay disappear.
\end{itemize}

\subsection{The Red Zone Protocol (Acute Crisis)}
Triggered manually (Panic Button) or automatically (Markov threshold breach).
\begin{itemize}
    \item \textbf{AI Triage:} Oloo shifts to a deterministic, grounding script. It does not use generative chat; it outputs a sequence of 5-4-3-2-1 grounding exercises.
    \item \textbf{Parallel Triage:} A modal pops up with two buttons:
    \begin{enumerate}
        \item \texttt{Call Befrienders Kenya} (Triggers native dialer).
        \item \texttt{Request Clinical Voucher} (Launches Triage Questionnaire).
    \end{enumerate}
    \item \textbf{Triage Questionnaire:} A rapid, validated assessment (e.g., GAD-7 or a custom gambling crisis scale). If the score crosses the clinical threshold, the API generates a Digital Voucher.
\end{itemize}

\section{Provider Portal (Therapist Dashboard)}
A secure web application for licensed clinicians. Accessible via \texttt{app.losslens.com/provider}.

\subsection{Provider Onboarding (Clinical Governance)}
\begin{itemize}
    \item \textbf{Credential Verification:} Therapists upload their Kenya Medical Practitioners and Dentists Council (KMPDC) license.
    \item \textbf{Admin Approval:} A LossLens Clinical Director (a licensed psychiatrist) reviews the license and approves the account.
    \item \textbf{Provider Profile:} Includes specialties (e.g., CBT, Gambling Addiction), languages spoken (English, Kiswahili, Sheng), and availability (Telehealth or Physical Clinic).
\end{itemize}

\subsection{Voucher Management Flow}
\begin{enumerate}
    \item \textbf{Voucher Queue:} The therapist sees a list of assigned Vouchers. A Voucher contains: \texttt{User\_ID}, \texttt{Voucher\_Code}, \texttt{Max\_Sessions} (e.g., 4), \texttt{Approved\_Fee\_per\_Session} (e.g., KES 3,000), and \texttt{Expiry\_Date} (e.g., 14 days).
    \item \textbf{Session Logging:} After a telehealth call or physical session, the therapist fills out the Shared Claim File.
    \item \textbf{Clinical Notes:} Free text field. Structured diagnosis code (ICD-11 for gambling disorder).
    \item \textbf{Submit Claim:} The therapist clicks submit. The system validates that the Voucher code is active and the \texttt{Max\_Sessions} limit has not been exceeded.
\end{enumerate}

\subsection{Payment \& Settlement Interface}
\begin{itemize}
    \item \textbf{Dashboard:} Shows "Pending Claims", "Approved Claims", and "Paid Claims".
    \item \textbf{Payout Gateway:} Integrates with M-Pesa B2C API (or a payment aggregator like Flutterwave). Once a claim is auto-approved, the payout is triggered within 48 hours.
    \item \textbf{Invoice Generation:} At the end of the month, the provider can download a consolidated PDF invoice of all paid claims for their own tax/accounting records.
\end{itemize}

\subsection{Immutable Audit Trail}
Every action (viewing a voucher, submitting notes, editing a note within a 15-minute window) is logged in the \texttt{audit\_logs} table. Therapists cannot delete claims once submitted; they can only request a void, which requires Admin approval.

% CONTINUATION OF LOSSLENS DOCUMENTATION
% Ensure you remove the \end{document} from Part 1 before appending this.

\section{Insurer Portal (B2B Dashboard)}
A secure web dashboard used by the insurance partner's claims adjusters, underwriters, and executives. Accessible via \texttt{app.losslens.com/insurer}.

\subsection{Population Health \& Actuarial Dashboard}
This is the primary view for the insurer. It relies entirely on aggregated, anonymized data. 
\begin{itemize}
    \item \textbf{Markov Transition Graphs:} Visualizes the movement of the insured cohort between the three states: \textit{Abstinent}, \textit{Lapsed}, and \textit{Relapsed}.
    \item \textbf{Claim Frequency Projections:} Based on the current Markov state distribution, the system projects the expected number of clinical claims for the next 30/60/90 days. This allows the insurer to manage their reserves.
    \item \textbf{Program ROI Tracking:} Compares the baseline expected claims (before LossLens) vs. actual claims. This is the core proof of the value proposition: \textit{"LossLens reduced your expected claim frequency by 22\% this quarter."}
\end{itemize}

\subsection{Voucher Pre-Authorization Flow}
LossLens eliminates retrospective denials by forcing pre-authorization.
\begin{enumerate}
    \item \textbf{The Trigger:} A user triggers a Clinical Voucher via the app's Triage Questionnaire.
    \item \textbf{The Notification:} The Insurer Portal receives a real-time webhook notification: \textit{New Voucher Requested: User \#402}.
    \item \textbf{The Review:} The insurer clicks the notification. They see the anonymized Triage Score and the proposed treatment (e.g., 4 Telehealth CBT sessions at KES 3,000 each).
    \item \textbf{The Approval:} The insurer clicks \texttt{Approve}. The Voucher is now locked and assigned to a therapist.
    \item \textbf{The Auto-Approve (Optional):} For Vouchers below a certain threshold (e.g., KES 5,000), the insurer can toggle "Auto-Approve" to bypass manual review entirely.
\end{enumerate}

\subsection{The Line-Item Ledger (Claims Adjudication)}
At the end of the month, LossLens submits a bulk invoice.
\begin{itemize}
    \item \textbf{The Ledger View:} The insurer sees an Excel/CSV export with every single session billed.
    \item \textbf{Columns:} \texttt{Claim\_ID}, \texttt{User\_Anonymized\_ID}, \texttt{Provider\_License}, \texttt{Session\_Date}, \texttt{ICD-11\_Code}, \texttt{Voucher\_Ref}, \texttt{Fee}.
    \item \textbf{The Audit Check:} The insurer can cross-reference the \texttt{Voucher\_Ref} (which they already approved) to ensure they are not being double-billed.
\end{itemize}

\subsection{Anti-Fraud Dashboard}
Tools to help the insurer trust the platform.
\begin{itemize}
    \item \textbf{Provider Outlier Detection:} Flags therapists whose billing frequency is statistically anomalous (e.g., a therapist billing for 20 crisis sessions in a week).
    \item \textbf{User Device Fingerprinting:} If the insurer suspects a specific claim is fraudulent, LossLens can verify that the user's device ID matches their historical login patterns.
    \item \textbf{Immutable Logs:} The insurer can view the full clinical note history for any claim. If a therapist tries to edit a note after a claim is paid, the original note remains visible in the audit log.
\end{itemize}

\section{Admin Portal (LossLens Internal)}
The master control panel for the LossLens team. Accessible via \texttt{app.losslens.com/admin}.

\subsection{Clinical Governance}
\begin{itemize}
    \item \textbf{Provider Management:} Approve new therapist registrations, suspend accounts for misconduct, or remove users from the network.
    \item \textbf{The Audit Tool:} A randomized list of 10\% of all paid claims is generated weekly. Admin staff must call the user to verify: \textit{"Did you actually speak to Dr. [Name] on [Date]?"} If the answer is "No," the claim is flagged as fraud and the therapist is investigated.
    \item \textbf{Clinical Director Oversight:} A dedicated dashboard for the LossLens Clinical Director to review flagged claims and ensure therapists are following LossLens CBT protocols.
\end{itemize}

\subsection{Financial Float Management}
LossLens pays therapists within 48 hours, relying on a float account before reimbursing from the insurer. The Admin Portal manages this.
\begin{itemize}
    \item \textbf{Float Balance Monitor:} A live view of the LossLens Bank/M-Pesa float balance.
    \item \textbf{Burn Rate Projection:} Predicts how many days until the float runs out based on current claim volume.
    \item \textbf{Reimbursement Tracker:} Tracks which insurers have paid their monthly bulk invoices and which are overdue. If an insurer is overdue by 30 days, the system automatically pauses new Voucher generation for that insurer's cohort.
\end{itemize}

\subsection{System Health \& Database Management}
\begin{itemize}
    \item \textbf{Supabase Connection Status:} A simple ping monitor to ensure the database is active and has not gone into the 7-day pause.
    \item \textbf{RLS Policy Verification:} A script that runs a daily audit to ensure Row Level Security is enabled on every table. If a table is found with RLS disabled, the Admin on duty is paged.
    \item \textbf{API Rate Limit Monitor:} Tracks usage against the Groq API limits to ensure the AI chat doesn't get throttled during peak usage.
\end{itemize}

\section{API \& Data Architecture}
The system relies on REST and webhook APIs to connect the React frontend, the Capacitor mobile app, and the Supabase backend.

\subsection{Core API Endpoints (Supabase Edge Functions)}
\begin{itemize}
    \item \texttt{POST /v1/triage/submit}: Accepts the Triage Questionnaire payload. Returns the Triage Score and (if approved) the Voucher Code.
    \item \texttt{POST /v1/markov/evaluate}: Accepts daily behavioral logs. Runs the transition matrix logic and returns the updated risk score.
    \item \texttt{GET /v1/provider/vouchers}: Returns a list of active vouchers for a specific provider ID.
    \item \texttt{POST /v1/provider/claims}: Submits a new claim for a specific voucher.
    \item \texttt{POST /v1/insurer/approve}: Insurer endpoint to approve a pending voucher.
    \item \texttt{POST /v1/webhooks/payment}: Listens for M-Pesa B2C success callbacks to update claim status.
\end{itemize}

\subsection{Database Schema (Core Tables)}
The Supabase PostgreSQL database is structured for high security and auditability.

\begin{lstlisting}[language=SQL]
-- Users Table
CREATE TABLE users (
    id UUID PRIMARY KEY DEFAULT uuid_generate_v4(),
    phone_number VARCHAR(15) UNIQUE NOT NULL,
    pgsi_score INTEGER,
    tier VARCHAR(10) DEFAULT 'FREE', -- 'FREE' or 'PREMIUM'
    shield_enabled BOOLEAN DEFAULT FALSE,
    created_at TIMESTAMP DEFAULT NOW()
);

-- Behavioral Logs (The Data Engine)
CREATE TABLE behavioral_logs (
    id UUID PRIMARY KEY,
    user_id UUID REFERENCES users(id),
    log_date DATE NOT NULL,
    mood_score INTEGER, -- 1-5
    urge_intensity INTEGER, -- 1-10
    gambled_today BOOLEAN,
    markov_risk_score DECIMAL(5,2),
    created_at TIMESTAMP DEFAULT NOW()
);

-- Vouchers (The Clinical Bridge)
CREATE TABLE vouchers (
    id UUID PRIMARY KEY,
    user_id UUID REFERENCES users(id),
    insurer_id UUID REFERENCES insurers(id),
    status VARCHAR(20) DEFAULT 'PENDING', -- PENDING, APPROVED, CLAIMED, EXPIRED
    max_sessions INTEGER DEFAULT 4,
    approved_fee_per_session DECIMAL(10,2),
    created_at TIMESTAMP DEFAULT NOW()
);

-- Claims (The Financial Record)
CREATE TABLE claims (
    id UUID PRIMARY KEY,
    voucher_id UUID REFERENCES vouchers(id),
    provider_id UUID REFERENCES providers(id),
    session_notes TEXT,
    claim_amount DECIMAL(10,2),
    status VARCHAR(20) DEFAULT 'SUBMITTED', -- SUBMITTED, PAID, REJECTED, FRAUD
    payout_mpesa_ref VARCHAR(50),
    created_at TIMESTAMP DEFAULT NOW()
);
\end{lstlisting}

\section{Edge Cases \& Failure Protocols}
A detailed technical document must account for what happens when things go wrong.

\subsection{Technical Failures}
\begin{itemize}
    \item \textbf{API Downtime (Groq/Supabase):} If the AI chat is down, the app falls back to a local database of pre-written CBT exercises. If Supabase is down, the app caches daily check-ins locally and pushes them to the server when connection is restored.
    \item \textbf{Payment Gateway Failure:} If the M-Pesa B2C API fails, the claim status is marked \texttt{PAYMENT\_FAILED}. The Admin Portal alerts the finance team, who trigger a manual payout within 24 hours.
    \item \textbf{Database Pause (7-day rule):} If the database pauses despite keep-alive scripts, the frontend shows a polite "We're doing maintenance" screen rather than an error stack.
\end{itemize}

\subsection{Clinical \& Fraud Failures}
\begin{itemize}
    \item \textbf{Provider Retraction:} If a therapist realizes a claim was submitted in error, they can click "Void Claim" within 15 minutes. After 15 minutes, they must email the Admin team to reverse the payout.
    \item \textbf{User Deletion Request:} If a user invokes their Right to Erasure (Kenya DPA 2019), their PII (phone number, name) is deleted immediately. However, their anonymized behavioral logs and financial claims are retained for 7 years to comply with insurance audit laws.
    \item \textbf{Insurer Rejection:} If an insurer rejects a claim post-payout, LossLens absorbs the loss and deducts it from the therapist's future payouts ONLY if fraud is proven. If the rejection is an administrative error by the insurer, LossLens disputes it using the immutable audit logs.
\end{itemize}

\section{Deployment \& DevOps}
\begin{itemize}
    \item \textbf{Web Frontend:} Deployed on Vercel. Continuous Integration via GitHub Actions. Every push to \texttt{main} triggers a production build.
    \item \textbf{Mobile App:} Built using Capacitor. Pushed to Google Play Store (Internal Testing Track) and Apple TestFlight.
    \item \textbf{Database:} Supabase. Backups are automated daily via Supabase Pro (once scaled) or manual pg\_dump scripts (during pilot).
\end{itemize}

% PART 3: EXISTING LOSSLENS FEATURES (ACCURATE TO CODEBASE)
% Insert after Executive Summary, before "The User Application (Module 1)"

\section{Core Platform Features (Existing Ecosystem)}
LossLens is not a future concept—it is a live, multi-tiered platform with distinct user journeys, community features, family support mechanisms, and B2B administration. This section documents the features that exist today.

\subsection{Public Access Layer}
\begin{itemize}
    \item \textbf{Landing Page:} A public-facing marketing and education page that introduces the platform, explains the PGSI framework, and directs users to sign up.
    \item \textbf{Pricing Page:} A transparent view of the Free vs. Premium tiers. Includes the micro-insurance upsell.
    \item \textbf{Family Join:} A unique invite link system that allows a family member to join the user's support network (see Family Portal below).
    \item \textbf{Authentication:} Signup, Login, Forgot Password, and Reset Password flows, secured via Supabase Auth.
\end{itemize}

\subsection{Onboarding \& Assessment}
\begin{itemize}
    \item \textbf{Guided Onboarding:} A multi-step flow that introduces the app, its features, and its philosophy of non-judgmental support.
    \item \textbf{PGSI Baseline Assessment:} A validated 9-question Problem Gambling Severity Index that determines the user's risk band (Low, Moderate, High).
    \item \textbf{Personalization Logic:} Based on the PGSI score, the platform adjusts the user's default dashboard, suggested features, and the frequency of support nudges.
\end{itemize}

\subsection{Daily Recovery Tools}
\begin{itemize}
    \item \textbf{Dashboard:} The central hub showing the user's risk score, daily check-in prompt, streak counter, and quick access to all tools.
    \item \textbf{Diary:} A private journaling space for free-text reflection, trigger logging, and mood tracking.
    \item \textbf{CBT Modules:} A structured library of Cognitive Behavioral Therapy lessons. Each module includes reading material, interactive exercises, and progress tracking.
    \item \textbf{Urge Toolkit:} A rapid-access page for in-the-moment crisis management. Includes guided breathing, the 5-4-3-2-1 grounding technique, and urge-surfing scripts.
\end{itemize}

\subsection{Community \& Social Features}
\begin{itemize}
    \item \textbf{The Circle:} A moderated community space where users can share their journey, post updates, and support one another. Includes guidelines enforcement and reporting tools.
    \item \textbf{Whispers:} An anonymous confession/vent channel. Users can post thoughts they are ashamed to attach to their identity. This lowers the barrier for disclosure of relapse.
    \item \textbf{Events:} A calendar of community events, recovery meetups, live therapy sessions, and webinars. Includes RSVP functionality.
\end{itemize}

\subsection{Commitment \& Reflection Tools}
\begin{itemize}
    \item \textbf{Precious Pact:} A commitment contract feature. Users publicly (or privately) pledge to a goal (e.g., "30 days clean"). Includes a "play count" tracker that logs how many times the user has revisited their commitment during an urge. The RPC \texttt{increment\_pact\_play\_count} suggests the app records these visits.
    \item \textbf{Monthly Mirror:} A guided audio/video recording tool where users record a message to their future self. If they relapse later, they can watch their own past self explain why they want to stay clean. Stored in the \texttt{mirror-recordings} Supabase bucket.
\end{itemize}

\subsection{Family \& Supporter Portal}
\begin{itemize}
    \item \textbf{Family Portal:} A dedicated view for a user's nominated family member or supporter. It provides aggregate wellness data (e.g., "Your family member checked in 5 times this week") without revealing private diary content or AI chats.
    \item \textbf{Family Join Link:} The user generates an invite link that gives a family member access to the Family Portal. This is opt-in and revocable at any time.
\end{itemize}

\subsection{Subscription \& Monetization}
\begin{itemize}
    \item \textbf{Subscription Page:} The self-service portal for upgrading to Premium. Displays the micro-insurance tier, benefits, and billing history.
    \item \textbf{Pricing Tier Enforcement:} The platform enforces feature gating (Free vs. Premium) at the route level via \texttt{ProtectedRoute}.
\end{itemize}

\subsection{AI Companion (Oloo)}
\begin{itemize}
    \item \textbf{24/7 Conversational Support:} An always-available AI companion accessed from the Dashboard.
    \item \textbf{CBT Integration:} Uses structured prompts to guide users through urge-surfing, reframing, and grounding.
    \item \textbf{Crisis Guardrails:} Detects crisis language and halts the generative response to present emergency hotlines (Befrienders Kenya).
\end{itemize}

\subsection{Resources \& Professional Referral}
\begin{itemize}
    \item \textbf{Resources Page:} A curated library of articles, videos, financial planning tools, and external links to professional help.
    \item \textbf{GAMHELP Pathway:} A dedicated referral flow that connects stabilized users to GAMHELP Kenya for long-term, in-person rehabilitation.
\end{itemize}

\section{Unified Admin Portal with Elevated Access}
All LossLens staff use a unified administrative portal with role-based elevated access. This is critical: LossLens is not just a consumer app—it is a B2B infrastructure play.

\subsection{Admin Dashboard}
All staff use the unified \texttt{/admin} portal. Administrative features are organized within this single portal rather than separate Super Admin routes.
\begin{itemize}
    \item \textbf{Admin Dashboard:} The operational hub for LossLens staff. Displays active user counts, crisis alerts, and platform health.
    \item \textbf{Manage Events:} Create, edit, and cancel community events.
    \item \textbf{Manage Groups:} Moderate The Circle community groups and enforce guidelines.
    \item \textbf{Manage Resources:} Curate the Resources library (add/remove articles and links).
    \item \textbf{User Analytics:} Aggregate statistics on user engagement, PGSI distribution, and relapse trends.
    \item \textbf{Broadcast:} Send platform-wide push notifications (e.g., "New CBT module launched").
    \item \textbf{Crisis Alerts:} A real-time dashboard that flags users whose Markov risk score has crossed a critical threshold. Admins can trigger manual outreach or dispatch clinical vouchers.
\end{itemize}

\subsection{Elevated Admin Features}
\begin{itemize}
    \item \textbf{User Management:} CRUD operations on all user accounts. Includes role assignment and account suspension.
    \item \textbf{Admin Management:} Create and revoke Admin accounts. Enforces separation of duties.
    \item \textbf{System Settings:} Toggle feature flags, adjust the Markov model thresholds, and configure pricing tiers.
    \item \textbf{Audit Logs:} An immutable, chronological log of every admin action. Critical for compliance and fraud prevention.
    \item \textbf{B2B Clients:} Manage institutional sponsors (SACCOs, employers, NGOs). Each B2B client can have their own cohort of users.
    \item \textbf{Insurance Partners:} Manage insurer integrations. Configure the pricing, voucher limits, and API endpoints for each insurance partner.
    \item \textbf{Grant Applications:} Track grant funding applications to NGOs, foundations, and government bodies. Displays status (Submitted, Approved, Rejected) and funding amounts.
    \item \textbf{Chat Viewer:} (Privacy-sensitive) A restricted tool that allows authorized administrators to review AI chat logs for quality assurance, crisis detection, and compliance audits. Access to this feature is logged in the Audit Logs.
\end{itemize}

\subsection{Step-Up Authentication Protocol}
\begin{itemize}
    \item All staff log in to a single \texttt{/admin} portal.
    \item Sensitive actions (User Management, Insurance Partners, B2B Clients, Chat Viewer, Audit Logs, System Settings) require the admin to re-enter their password.
    \item Upon password entry, elevated permissions unlock for exactly 15 minutes.
    \item Every unlock is logged in the \texttt{audit\_logs} table.
\end{itemize}

\subsection{Database-Level Enforcement}
Row Level Security (RLS) is used to hide sensitive tables, such as \texttt{insurance\_partners}, from normal administrators. Elevated access is enforced by database policies rather than relying only on frontend route controls.

\lstset{language=SQL}
\begin{lstlisting}
CREATE POLICY "Super admins can view insurance partners"
ON public.insurance_partners FOR SELECT
USING (
    EXISTS (
        SELECT 1 FROM public.profiles
        WHERE profiles.id = auth.uid()
        AND profiles.is_super_admin = TRUE
    )
);
\end{lstlisting}

\subsection{The Immutable Audit Log}
The \texttt{audit\_logs} table allows INSERT and SELECT operations for super administrators, but it has NO UPDATE or DELETE policies. This cryptographically immutable audit trail prevents authorized and unauthorized users from modifying or removing audit records after they are created.

\section{Data \& Privacy Architecture (Existing)}
The existing platform already enforces several important data separation principles.

\subsection{Row Level Security (RLS)}
All Supabase tables are protected by RLS policies. Users can only access their own rows. Admins can access aggregate data. Super Admins have scoped access governed by the Audit Logs.

\subsection{Storage Buckets}
The platform uses Supabase Storage with separate buckets for:
\begin{itemize}
    \item \texttt{mirror-recordings}: Encrypted audio/video files for the Monthly Mirror feature.
    \item \texttt{precious-pacts}: Commitment contract documents and evidence.
\end{itemize}

\subsection{The Privacy Firewall Principle}
Even in the current implementation, the platform follows the principle that sensitive recovery data (diary, AI chat, PGSI) should not be exposed to B2B partners. The Insurer Portal (planned) will only see anonymized, aggregate Markov output.

\section{Appendix A: Master Database Migration Script}
The following master migration script is used to deploy the LossLens database infrastructure.

\lstset{language=SQL}
\begin{lstlisting}
-- LOSSLENS MASTER MIGRATION SCRIPT
-- 1. Add Super Admin flag to profiles
ALTER TABLE public.profiles
ADD COLUMN IF NOT EXISTS is_super_admin BOOLEAN DEFAULT FALSE;

-- 2. Create Audit Logs table
CREATE TABLE IF NOT EXISTS public.audit_logs (
    id UUID PRIMARY KEY DEFAULT gen_random_uuid(),
    admin_id UUID REFERENCES public.profiles(id),
    action TEXT NOT NULL,
    target_table TEXT,
    target_id UUID,
    metadata JSONB,
    created_at TIMESTAMPTZ DEFAULT NOW()
);

-- 3. Enable RLS on sensitive tables
ALTER TABLE public.profiles ENABLE ROW LEVEL SECURITY;
ALTER TABLE public.audit_logs ENABLE ROW LEVEL SECURITY;

-- 4. Create RLS Policies
DROP POLICY IF EXISTS "Super admins can read audit logs" ON public.audit_logs;

CREATE POLICY "Super admins can read audit logs"
ON public.audit_logs FOR SELECT
USING (
    EXISTS (
        SELECT 1 FROM public.profiles
        WHERE profiles.id = auth.uid()
        AND profiles.is_super_admin = TRUE
    )
);

CREATE POLICY "Admins can insert audit logs"
ON public.audit_logs FOR INSERT
WITH CHECK (
    EXISTS (
        SELECT 1 FROM public.profiles
        WHERE profiles.id = auth.uid()
        AND profiles.role = 'admin' OR profiles.is_super_admin = TRUE
    )
);

-- 5. Reload Schema
NOTIFY pgrst, 'reload schema';
\end{lstlisting}

\end{document}